\documentclass[a4paper]{arrowhead}

\usepackage[yyyymmdd]{datetime}
\usepackage{etoolbox}
\usepackage[utf8]{inputenc}
\usepackage[T1]{fontenc}
\usepackage{multirow}
\usepackage{longtable}
\usepackage{tikz}
\usetikzlibrary{positioning,shapes.multipart,arrows.meta}

\renewcommand{\dateseparator}{-}
\setlength{\emergencystretch}{3em}
\lstset{language=HTTP,breakatwhitespace=false,columns=fullflexible}

%% Special references
\newcommand{\mref}[2][]{{\textcolor{ArrowheadPurple}{\hyperref[sec:model:#2]{\ifstrempty{#1}{#2}{#1}}}}}
\newcommand{\oref}[1]{{\textcolor{ArrowheadBlue}{\hyperref[sec:operations:#1]{#1}}}}
\newcommand{\pdef}[1]{{\textcolor{ArrowheadGrey}{#1\label{sec:model:primitives:#1}\label{sec:model:primitives:#1s}\label{sec:model:primitives:#1es}}}}
\newcommand{\pref}[1]{{\textcolor{ArrowheadGrey}{\hyperref[sec:model:primitives:#1]{#1}}}}

% as template_sd.tex; the input and output arguments are passed already formatted (\mref or \pref)
\newrobustcmd\osubsection[3]{
  \addtocounter{subsection}{1}
  \addcontentsline{toc}{subsection}{\protect\numberline{\thesubsection}operation \textcolor{ArrowheadBlue}{#1}}
  \subsection*{
    \thesubsection\quad
    operation
    \textcolor{ArrowheadBlue}{#1}
    (\notblank{#2}{#2}{})
    \notblank{#3}{: #3}{}
  }
  \label{sec:operations:#1}
}
\newrobustcmd\msubsection[2]{
  \addtocounter{subsection}{1}
  \addcontentsline{toc}{subsection}{\protect\numberline{\thesubsection}#1 \textcolor{ArrowheadPurple}{#2}}
  \subsection*{\thesubsection\quad#1 \textcolor{ArrowheadPurple}{#2}}
  \label{sec:model:#2} \label{sec:model:#2s} \label{sec:model:#2es}
}
%%

\begin{document}
\sloppy

%% Arrowhead Document Properties
\ArrowheadTitle{sysMonitorData}
\ArrowheadServiceID{sysMonitorData}
\ArrowheadType{Service Description}
\ArrowheadTypeShort{SD}
\ArrowheadVersion{5.2.2}
\ArrowheadDate{2026-10-30}
\ArrowheadAuthor{Jerker Delsing}
\ArrowheadStatus{RELEASE CANDIDATE}
\ArrowheadContact{jerker.delsing@ltu.se}
\ArrowheadFooter{\href{www.arrowhead.eu}{www.arrowhead.eu}}
\ArrowheadSetup
%%

%% Front Page
\begin{center}
  \vspace*{1cm}
  \huge{\arrowtitle}

  \vspace*{0.2cm}
  \LARGE{\arrowtype}
  \vspace*{1cm}

  \Large{Service ID: \textit{"\arrowid"}}

  \vspace*{\fill}

  % Front Page Image

  \vspace*{1cm}
  \vspace*{\fill}

  % Front Page Abstract
  \begin{abstract}
    This document describes the \texttt{sysMonitorData} service, version 1.0.0, produced by the SysMonitor support system of Arrowhead 5.2.2.
  \end{abstract}

  \vspace*{1cm}
\end{center}

\newpage
%%

%% Table of Contents
\tableofcontents
\newpage
%%

\section{Overview}
\label{sec:overview}

This document describes the \texttt{sysMonitorData} service, which provides the stream of all received observations, republished on the local MQTT broker.
It is produced by the SysMonitor system \cite{smo-sysd} in service version 1.0.0.

The rest of this document is organized as follows.
In Section \ref{sec:operations}, we describe the abstract message operations provided by the service.
In Section \ref{sec:model}, we end the document by presenting the data types used by the mentioned operations.
The concrete interface is specified in the MQTT/TLS/JSON IDD \cite{smo-sysMonitorData-mqtt-idd}.

\subsection{Significant Prior Art}

New in this workspace (SysMonitor-Data-MQTT-IDD-5.2.2).

\subsection{How This Service Is Meant to Be Used}

Consumers such as AIGym orchestrate \texttt{sysMonitorData} and subscribe to \texttt{arrowhead/\allowbreak{}\textless{}Cloud\textgreater{}/\allowbreak{}sysmonitor/\allowbreak{}data/\allowbreak{}\textless{}observation\allowbreak{}Type\textgreater{}/\allowbreak{}\textless{}provider\allowbreak{}System\textgreater{}}.

\subsection{Important Delimitations and Dependencies}
\label{sec:delimitations}

QoS 1 from a transactional outbox; duplicates are possible after a crash between PUBACK and outbox deletion. Broker ACLs decide who may subscribe.

\newpage

\section{Service Interface}
\label{sec:operations}

This section describes the interfaces to the \texttt{sysMonitorData} service.
In particular, each subsection names an abstract operation, an input type and an output type, in that order.
The input type is named inside parentheses, while the output type is preceded by a colon.
Input and output types are only denoted when accepted or returned, respectively, by the interface in question.

All abstract data types named in this section are defined in Section \ref{sec:model}.

\begin{figure}[ht!]
  \centering
  \begin{tikzpicture}[font=\small]
    \node[draw=ArrowheadBlue, very thick, rectangle split, rectangle split parts=2, align=center,
          rectangle split part fill={ArrowheadBlue!12,white}, text width=9cm, inner sep=6pt] (b)
      {\guillemotleft{}Service\guillemotright{}\\\textbf{sysMonitorData}
        \nodepart{second} \texttt{subscribe}};
  \end{tikzpicture}
  \caption{SysML-style block of the \texttt{sysMonitorData} service and its operations.}
  \label{fig:SD}
\end{figure}

The following interface operations are available.

\osubsection{subscribe}{}{\mref{StoredObservation}}



\newpage

\section{Information Model}
\label{sec:model}

Here, all data objects that can be part of the \texttt{sysMonitorData} service are listed in alphabetic order.
Note that each subsection, which describes one type of object, begins with the \textit{struct} keyword, which is used to denote a collection of named fields, each with its own data type.
As a complement to the explicitly defined types in this section, there is also a list of implicit primitive types in Section \ref{sec:model:primitives}, which are used to represent things like hashes and identifiers.

\msubsection{struct}{StoredObservation}




\noindent\begin{tabularx}{\textwidth}{| >{\raggedright\arraybackslash}p{3.9cm} | >{\raggedright\arraybackslash\small}p{5.3cm} | X |} \hline
\rowcolor{gray!33} Field & Type & Description \\ \hline
schema\allowbreak{}Version & \pref{String} & \texttt{1.\allowbreak{}0}. \\ \hline
observation\allowbreak{}Id & \pref{Identifier} & Identifier (publisher's, or a SHA-256 of routing key and content). \\ \hline
provider\allowbreak{}System & \pref{String} & Publishing system. \\ \hline
service\allowbreak{}Instance & \pref{String} & Service instance. \\ \hline
observation\allowbreak{}Type & \pref{String} & Observation type. \\ \hline
source\allowbreak{}Time & \pref{DateTime} & Publisher time, UTC. \\ \hline
received\allowbreak{}Time & \pref{DateTime} & Time SysMonitor received it. \\ \hline
published\allowbreak{}Time & \pref{DateTime} & Time it was forwarded on \texttt{sysMonitorData}. \\ \hline
value & \pref{Map}$<$\allowbreak{}\pref{String}, \pref{Object}$>$ & Observed value. \\ \hline
unit & \pref{String} & Unit. \\ \hline
quality & \pref{String} & Quality. \\ \hline
pipeline & \pref{Map}$<$\allowbreak{}\pref{String}, \pref{Object}$>$ & Intake details (routing key, exchange). \\ \hline
\end{tabularx}



\subsection{Primitives}
\label{sec:model:primitives}

Types and structures mentioned throughout this document that are assumed to be available to implementations of this service.
The concrete interpretations of each of these types and structures must be provided by any IDD document claiming to implement this service.

\noindent\begin{tabularx}{\textwidth}{| p{4.6cm} | X |} \hline
\rowcolor{gray!33} Type & Description \\ \hline
\pdef{DateTime} & Pinpoints a specific moment in time (UTC). \\ \hline
\pdef{Identifier} & A string that identifies a record, e.g. a UUID or a composite identifier. \\ \hline
\pdef{Map} & A set of key-value pairs. \\ \hline
\pdef{Object} & Any structured value. \\ \hline
\pdef{String} & A chain of characters. \\ \hline
\end{tabularx}

\newpage

\bibliographystyle{IEEEtran}
\bibliography{bibliography}

\newpage

\section{Revision History}
\subsection{Amendments}

\noindent\begin{tabularx}{\textwidth}{| p{1cm} | p{3cm} | p{2cm} | X | p{4cm} |} \hline
\rowcolor{gray!33} No. & Date & Version & Subject of Amendments & Author \\ \hline

1 & 2026-10-22 & \arrowversion & Document created from the Arrowhead LaTeX templates for the 5.2.2 Java Spring implementation, including the local extensions of this workspace. & Jerker Delsing \\ \hline
2 & 2026-10-30 & \arrowversion & Checked against the primary Java source tree (Eclipse-Arrowhead-java): services, operations, paths and data models unchanged. & Claude Opus 5.5 (draft) \\ \hline
\end{tabularx}

\subsection{Quality Assurance}

\noindent\begin{tabularx}{\textwidth}{| p{1cm} | p{3cm} | p{2cm} | X |} \hline
\rowcolor{gray!33} No. & Date & Version & Approved by \\ \hline

1 & & \arrowversion  &  \\ \hline

\end{tabularx}

\end{document}
